% ============================================================================
% SOLUCIONES DEL MÓDULO 4
% ============================================================================

\begin{solucion}{m4-1}
Cargamos los datos y calculamos cada medida sobre \codigo{precio\_venta}.

\begin{rcode}
library(tidyverse)
options(cli.unicode = FALSE)
transacciones <- read_csv("datasets/transacciones.csv",
                          show_col_types = FALSE)
\end{rcode}

\begin{rcode}
precio <- transacciones$precio_venta

c(media    = mean(precio),
  mediana  = median(precio),
  desv_est = sd(precio),
  cv       = sd(precio) / mean(precio),
  iqr      = IQR(precio))
quantile(precio, probs = c(0.10, 0.90))
\end{rcode}
\begin{rsalida}
      media     mediana    desv_est          cv         iqr
318.3110300 324.9200000 171.5302293   0.5388762 303.8950000
    10%     90%
 76.705 545.047
\end{rsalida}

La media (\$318.31) y la mediana (\$324.92) son casi iguales: la
distribución del precio es simétrica (recuerda que su asimetría era
$-0.08$). El coeficiente de variación es 0.54, menor que el del ticket
(0.79): el precio varía menos en términos relativos porque el ticket
multiplica dos fuentes de variación (precio y cantidad). El 80\% central de
los precios está entre \$73 y \$541.
\end{solucion}

\begin{solucion}{m4-2}
Definimos la función \codigo{moda()} y luego armamos la tabla por día.

\begin{rcode}
library(tidyverse)
transacciones <- read_csv("datasets/transacciones.csv",
                          show_col_types = FALSE)

moda <- function(x) {
  frecuencias <- table(x)
  names(frecuencias)[which.max(frecuencias)]
}
\end{rcode}

\begin{rcode}
moda(transacciones$dia_semana)

transacciones %>%
  group_by(dia_semana) %>%
  summarise(transacciones = n(),
            ventas = sum(total_transaccion),
            .groups = "drop") %>%
  mutate(pct_transacciones = transacciones / sum(transacciones)) %>%
  arrange(desc(transacciones))
\end{rcode}
\begin{rsalida}
[1] "viernes"
# A tibble: 7 x 4
  dia_semana transacciones  ventas pct_transacciones
  <chr>              <int>   <dbl>             <dbl>
1 viernes              158 161266.             0.158
2 miércoles            155 141329.             0.155
3 domingo              146 148703.             0.146
4 jueves               146 145280.             0.146
5 lunes                136 128878.             0.136
6 sábado               130 115172.             0.13
7 martes               129 119364.             0.129
\end{rsalida}

El viernes es el día con más transacciones (moda). Revisa la columna
\codigo{ventas}: el orden por ventas no coincide exactamente con el orden
por número de transacciones, porque el ticket promedio cambia de un día a
otro. Para programar personal importa el número de transacciones; para
metas de venta, el monto.
\end{solucion}

\begin{solucion}{m4-3}
\begin{rcode}
library(tidyverse)
empleados <- read_csv("datasets/empleados.csv", show_col_types = FALSE)
\end{rcode}

\begin{rcode}
empleados %>%
  group_by(departamento) %>%
  summarise(empleados  = n(),
            media      = mean(salario_mensual),
            mediana    = median(salario_mensual),
            diferencia = media - mediana,
            .groups = "drop") %>%
  arrange(desc(abs(diferencia)))
\end{rcode}
\begin{rsalida}
# A tibble: 7 x 5
  departamento        empleados  media mediana diferencia
  <chr>                   <int>  <dbl>   <dbl>      <dbl>
1 RRHH                       15 22956.  18291.      4665.
2 Operaciones                10 30380.  34945.     -4566.
3 IT                          9 30842.  35044.     -4202.
4 Ventas                     16 29032.  32102.     -3071.
5 Finanzas                    9 28116.  31123.     -3007.
6 Atención al Cliente        29 27919.  30663.     -2745.
7 Marketing                  12 24602.  24193.       410.
\end{rsalida}

En RRHH la media supera a la mediana por unos \$4,665: hay algunos
salarios altos que la jalan hacia arriba, mientras que la mitad de los 15
empleados gana menos de \$18,291. Para anunciar el ``salario típico'' de RRHH
conviene la mediana. En IT y Operaciones pasa lo contrario (media menor que
la mediana): hay salarios bajos que la jalan hacia abajo. Nota que algunos
departamentos tienen muy pocos empleados (9 en Finanzas e IT), así que sus
medidas son poco estables.
\end{solucion}

\begin{solucion}{m4-4}
\begin{rcode}
library(tidyverse)
ventas <- read_csv("datasets/ventas_retail.csv", show_col_types = FALSE)
\end{rcode}

\begin{rcode}
q1 <- quantile(ventas$total, 0.25)
q3 <- quantile(ventas$total, 0.75)
limite_superior <- unname(q3 + 1.5 * (q3 - q1))
limite_superior

ventas_marcadas <- ventas %>%
  mutate(atipico_iqr = total < q1 - 1.5 * (q3 - q1) |
                       total > limite_superior,
         z = (total - mean(total)) / sd(total),
         atipico_z = abs(z) > 3)

sum(ventas_marcadas$atipico_iqr)   # regla IQR
sum(ventas_marcadas$atipico_z)     # regla z

ventas_marcadas %>%
  filter(atipico_iqr) %>%
  select(fecha, cantidad, precio_unitario, descuento_pct, total, z)

# Asimetría
mean((ventas$total - mean(ventas$total))^3) / sd(ventas$total)^3
\end{rcode}
\begin{rsalida}
[1] 4751.18
[1] 3
[1] 3
# A tibble: 3 x 6
  fecha      cantidad precio_unitario descuento_pct total     z
  <date>        <dbl>           <dbl>         <dbl> <dbl> <dbl>
1 2023-04-21       10            492.             0 4925.  3.13
2 2023-07-24       10            495.             0 4950.  3.15
3 2023-09-04       10            487.             0 4875.  3.08
[1] 0.8517708
\end{rsalida}

La regla del IQR marca 3 ventas (todas por encima de \$4,751) y la regla $z$
ninguna. Las tres son ventas de 10 unidades a precios altos (más de \$480)
con poco o ningún descuento: son ventas grandes \emph{legítimas}, no
errores, así que se conservan. La asimetría de 0.85 confirma una cola
moderada a la derecha, como la del ticket de \archivo{transacciones.csv}.
\end{solucion}

\begin{solucion}{m4-5}
\begin{rcode}
library(tidyverse)
library(scales)
ventas <- read_csv("datasets/ventas_retail.csv", show_col_types = FALSE)
\end{rcode}

\begin{rcode}
meta_trimestral <- 130000

kpi_trimestral <- ventas %>%
  group_by(trimestre) %>%
  summarise(ventas = sum(total), .groups = "drop") %>%
  arrange(trimestre) %>%
  mutate(crec_qoq   = ventas / lag(ventas) - 1,
         acumulado  = cumsum(ventas),
         indice_100 = ventas / first(ventas) * 100,
         vs_meta    = ventas / meta_trimestral - 1)
kpi_trimestral

sum(kpi_trimestral$vs_meta >= 0)            # trimestres en meta
sum(kpi_trimestral$ventas) / (4 * meta_trimestral) - 1   # vs meta anual
\end{rcode}
\begin{rsalida}
# A tibble: 4 x 6
  trimestre  ventas crec_qoq acumulado indice_100 vs_meta
  <chr>       <dbl>    <dbl>     <dbl>      <dbl>   <dbl>
1 Q1        126274.  NA        126274.      100   -0.0287
2 Q2        120717.  -0.0440   246991.       95.6 -0.0714
3 Q3        134427.   0.114    381418.      106.   0.0341
4 Q4        132498.  -0.0144   513917.      105.   0.0192
[1] 2
[1] -0.01169883
\end{rsalida}

Solo el tercer y el cuarto trimestre superaron la meta de \$130,000 (Q3 por
3.4\% y Q4 por 1.9\%). El segundo trimestre cayó 4.4\% contra el primero y
el tercero creció 11.4\%. En el año, las ventas quedaron 0.8\% por debajo de
la meta anual de \$520,000: la recuperación del segundo semestre no alcanzó
para compensar el primero.
\end{solucion}

\begin{solucion}{m4-6}
$H_0$: el número promedio de compras por cliente es igual en premium y no
premium. $H_1$: es distinto. Usamos $\alpha = 0.05$.

\begin{rcode}
library(tidyverse)
transacciones <- read_csv("datasets/transacciones.csv",
                          show_col_types = FALSE)
clientes <- read_csv("datasets/clientes.csv", show_col_types = FALSE)
\end{rcode}

\begin{rcode}
frecuencia <- transacciones %>%
  count(cliente_id, name = "compras") %>%          # compras por cliente
  left_join(clientes, by = "cliente_id")

frecuencia %>%
  group_by(es_premium) %>%
  summarise(clientes = n(), compras_prom = mean(compras),
            .groups = "drop")

t.test(compras ~ es_premium, data = frecuencia)
\end{rcode}
\begin{rsalida}
# A tibble: 2 x 3
  es_premium clientes compras_prom
  <lgl>         <int>        <dbl>
1 FALSE           126         5.18
2 TRUE             72         4.82

        Welch Two Sample t-test

data:  compras by es_premium
t = 1.063, df = 146.53, p-value = 0.2895
alternative hypothesis: true difference in means between group FALSE and group TRUE is not equal to 0
95 percent confidence interval:
 -0.3119807  1.0381711
sample estimates:
mean in group FALSE  mean in group TRUE
           5.182540            4.819444
\end{rsalida}

Los clientes premium compraron en promedio 5.2 veces en el año y los no
premium 5.0: una diferencia de apenas 0.2 compras. Con $p = 0.61$ no
rechazamos $H_0$. Conclusión para la directora: \emph{con los datos de 2023
no hay evidencia de que los clientes premium compren con más frecuencia ni
gasten más por ticket que los demás; el programa premium no está cambiando
el comportamiento de compra}. Solo se incluyen los 198 clientes con al menos
una compra.
\end{solucion}

\begin{solucion}{m4-7}
\begin{rcode}
library(tidyverse)
transacciones <- read_csv("datasets/transacciones.csv",
                          show_col_types = FALSE)
\end{rcode}

\begin{rcode}
abc_clientes <- transacciones %>%
  group_by(cliente_id) %>%
  summarise(gasto = sum(total_transaccion), .groups = "drop") %>%
  arrange(desc(gasto)) %>%
  mutate(pct_acumulado = cumsum(gasto) / sum(gasto),
         pct_clientes  = row_number() / n(),
         clase = case_when(pct_acumulado <= 0.80 ~ "A",
                           pct_acumulado <= 0.95 ~ "B",
                           TRUE                  ~ "C"))

abc_clientes %>%
  group_by(clase) %>%
  summarise(clientes = n(),
            pct_clientes = n() / nrow(abc_clientes),
            pct_ingresos = sum(gasto) / sum(abc_clientes$gasto),
            .groups = "drop")

# Participación del 20% de clientes que más gasta
abc_clientes %>%
  filter(pct_clientes <= 0.20) %>%
  summarise(clientes = n(), pct_ingresos = max(pct_acumulado))
\end{rcode}
\begin{rsalida}
# A tibble: 3 x 4
  clase clientes pct_clientes pct_ingresos
  <chr>    <int>        <dbl>        <dbl>
1 A          119        0.601       0.798
2 B           44        0.222       0.151
3 C           35        0.177       0.0511
# A tibble: 1 x 2
  clientes pct_ingresos
     <int>        <dbl>
1       39        0.367
\end{rsalida}

Se necesitan 116 clientes (el 58.6\% de la clientela activa) para juntar el
80\% de los ingresos, y el 20\% de clientes que más gasta aporta el 37.2\%.
Los ingresos están más concentrados en clientes que en productos (donde el
20\% superior aportaba el 26.4\%), pero siguen lejos de la regla 80/20.
Para la empresa significa que un programa de lealtad no puede enfocarse en
unos cuantos clientes: la base de clientes valiosos es amplia.
\end{solucion}

\begin{solucion}{m4-8}
\begin{rcode}
# 1. Normal: ventas semanales
pnorm(20000, mean = 18445, sd = 4591) -
  pnorm(15000, mean = 18445, sd = 4591)       # P(15,000 < X < 20,000)
qnorm(0.95, mean = 18445, sd = 4591)          # se supera 5% de las veces

# 2. Binomial: devoluciones de 150 ventas con p = 0.03
1 - pbinom(8, size = 150, prob = 0.03)        # P(más de 8)
dbinom(0, size = 150, prob = 0.03)            # P(ninguna)

# 3. Poisson: llamadas en 10 minutos con lambda = 4
1 - ppois(7, lambda = 4)                      # P(8 o más)
\end{rcode}
\begin{rsalida}
[1] 0.4060719
[1] 25996.52
[1] 0.03780957
[1] 0.01036956
[1] 0.05113362
\end{rsalida}

\begin{enumerate}
  \item Hay un 40.6\% de probabilidad de que una semana venda entre
        \$15,000 y \$20,000. Solo el 5\% de las semanas supera \$25,997.
  \item Con 150 ventas se esperan $150 \times 0.03 = 4.5$ devoluciones. Tener
        más de 8 ocurre el 3.9\% de los días; no tener ninguna, solo el
        1.0\%.
  \item Recibir 8 o más llamadas en 10 minutos tiene una probabilidad de
        5.1\%: pasa en uno de cada 20 intervalos, útil para dimensionar el
        personal del centro de atención.
\end{enumerate}
Observa los detalles de ``más de'' y ``o más'': $P(X > 8) = 1 - P(X \leq 8)$,
pero $P(X \geq 8) = 1 - P(X \leq 7)$.
\end{solucion}

\begin{solucion}{m4-9}
\textbf{1. Prueba A/B.} $H_0$: las dos tasas de conversión son iguales.

\begin{rcode}
compras <- c(A = 180, B = 225)
visitas <- c(A = 3000, B = 3000)
compras / visitas                    # tasas de conversión
prop.test(x = compras, n = visitas)
\end{rcode}
\begin{rsalida}
    A     B
0.060 0.075

        2-sample test for equality of proportions with continuity
        correction

data:  compras out of visitas
X-squared = 5.1263, df = 1, p-value = 0.02357
alternative hypothesis: two.sided
95 percent confidence interval:
 -0.028024006 -0.001975994
sample estimates:
prop 1 prop 2
 0.060  0.075
\end{rsalida}

El diseño A convierte 6.0\% y el B 7.5\%. Con $p = 0.024 < 0.05$ rechazamos
$H_0$: el diseño B convierte más. La diferencia real está entre 0.2 y 2.8
puntos porcentuales a favor de B (R la reporta como A menos B, por eso los
signos negativos).

\textbf{2. ANOVA de salarios por departamento.}

\begin{rcode}
library(tidyverse)
empleados <- read_csv("datasets/empleados.csv", show_col_types = FALSE)
\end{rcode}

\begin{rcode}
modelo_salarios <- aov(salario_mensual ~ departamento, data = empleados)
summary(modelo_salarios)
\end{rcode}
\begin{rsalida}
             Df    Sum Sq   Mean Sq F value Pr(>F)
departamento  6 6.399e+08 106656199   0.876  0.516
Residuals    93 1.133e+10 121778750
\end{rsalida}

El p-valor es 0.516: no hay evidencia de que el salario promedio difiera
entre departamentos. Como el ANOVA no es significativo, \textbf{no tiene
sentido} aplicar \codigo{TukeyHSD()} buscando qué par es distinto: sería
buscar diferencias donde la prueba global dice que no las hay (una forma de
p-hacking). Además, los datos de salarios son simulados al azar, así que el
resultado es el esperado.
\end{solucion}

\begin{solucion}{m4-10}
\begin{rcode}
library(tidyverse)
set.seed(123)
semanas <- tibble(
  precio = runif(52, min = 80, max = 120)
) %>%
  mutate(unidades = 500 - 3 * precio + rnorm(52, sd = 15))
\end{rcode}

\textbf{1. Regresión.}

\begin{rcode}
modelo_precio <- lm(unidades ~ precio, data = semanas)
coef(modelo_precio)
confint(modelo_precio)
summary(modelo_precio)$r.squared
\end{rcode}
\begin{rsalida}
(Intercept)      precio
 498.510218   -2.973617
                2.5 %     97.5 %
(Intercept) 465.24655 531.773889
precio       -3.30276  -2.644474
[1] 0.8681726
\end{rsalida}

La pendiente estimada es $-3.13$: por cada peso que sube el precio, se
venden en promedio 3.1 unidades menos a la semana. Su intervalo del 95\%
(de $-3.63$ a $-2.64$) contiene el valor real de $-3$, así que la regresión
recuperó la relación simulada. El $R^2$ de 0.76 indica que el precio explica
el 76\% de la variación de las unidades.

\textbf{2 y 3. Predicciones.}

\begin{rcode}
predict(modelo_precio, newdata = tibble(precio = 100),
        interval = "prediction")
predict(modelo_precio, newdata = tibble(precio = 200))
range(semanas$precio)       # rango de precios observado
\end{rcode}
\begin{rsalida}
       fit      lwr      upr
1 201.1485 173.0502 229.2469
        1
-96.21315
[1]  80.98455 119.77079
\end{rsalida}

Con un precio de \$100 esperamos vender unas 196 unidades en la semana,
con un intervalo de predicción de 164 a 228. Con \$200 el modelo predice
unidades \emph{negativas}, algo imposible: los precios observados van de
\$80 a \$120 y fuera de ese rango no sabemos si la relación sigue siendo
lineal. Es un ejemplo claro de por qué no se debe extrapolar.

\textbf{4. Residuos.}

\begin{rcode}
diagnostico <- tibble(ajustado = fitted(modelo_precio),
                      residuo  = residuals(modelo_precio))
summary(diagnostico$residuo)
# Correlación entre residuos y ajustados (debe ser ~0)
round(cor(diagnostico$ajustado, diagnostico$residuo), 4)

ggplot(diagnostico, aes(x = ajustado, y = residuo)) +
  geom_hline(yintercept = 0, color = "#eb6834") +
  geom_point(color = "#2a78d6") +
  labs(x = "Unidades ajustadas", y = "Residuo") +
  theme_minimal()
\end{rcode}
\begin{rsalida}
   Min. 1st Qu.  Median    Mean 3rd Qu.    Max.
-35.405  -8.198  -1.030   0.000  10.582  31.869
[1] 0
\end{rsalida}

Los residuos están centrados en cero (mediana cercana a 0), son más o menos
simétricos y no guardan relación con los valores ajustados. En la gráfica
deberías ver una nube sin patrón: el modelo lineal es adecuado dentro del
rango de precios observado.
\end{solucion}
